\documentclass[10pt,a4paper]{amsart}
\usepackage[margin=19mm]{geometry}
\usepackage{amsmath,amssymb,mathtools}
\usepackage[scr=rsfs,cal=euler]{mathalfa}
\usepackage{xfrac}
\usepackage[unicode,hidelinks,bookmarks=false]{hyperref}
\newcommand{\ie}{i.e.\ }
\newcommand{\cf}{cf.\ }
\newcommand{\resp}{respectively\ }
\newcommand{\Subsection}{\subsection}
\providecommand{\nobreakdash}{\nobreak\hbox{-}}
\setlength{\emergencystretch}{2em}
\multlinegap0pt
\usepackage{bm}
\usepackage{bbm}
\usepackage{dsfont}
\usepackage{xspace}
\usepackage{cancel}
\usepackage{mathtools}
\usepackage{amsthm}
\makeatletter
\@ifundefined{theorem}{\newtheorem{theorem}{Theorem}}{}
\@ifundefined{lemma}{\newtheorem{lemma}[theorem]{Lemma}}{}
\makeatother
\DeclareMathOperator{\Eff}{\overline{Eff}}
\DeclareMathOperator{\NS}{NS}
\DeclareMathOperator{\adj}{adj}
\newcommand{\A}{\mathbb{A}}
\newcommand{\C}{\mathbb{C}}
\newcommand{\R}{\mathbb{R}}
\renewcommand{\Re}{\mathrm{Re}}
\title[Manin's conjecture for integral points on toric varieties]{Manin's conjecture for integral points on toric varieties}
\author{Tim Santens}
\date{Source: arXiv:2312.13914v2 (CC BY 4.0)}
\begin{document}
\maketitle

\noindent\emph{Source and licence.} Manin's conjecture for integral points on toric varieties. Version arXiv:2312.13914v2 (CC BY 4.0), distributed under the Creative Commons Attribution 4.0 International licence, as recorded in the arXiv licence metadata for this exact version and shown on the abstract page https://arxiv.org/abs/2312.13914v2. Full rights notice: licenses/arxiv-2312.13914-CC-BY-4.0.txt.\par\smallskip

\noindent\textit{Editorial scope.} Counted unit: the Lemma whose source label is \texttt{lem: Moving the contours} in the article (source0.tex lines 2453--2476, source label \texttt{lem: Moving the contours}), delivered together with the article's own complete original proof of it (source0.tex lines 2477--2533, one \texttt{proof} environment of 57 lines). No mathematics is written, repaired or re-derived: statement and proof are the article's own text line for line. Required context retained uncounted: Poles of characteristic functions of cones (lines 442--448); eq:log-canonical divisor (lines 2077--2079); lem:Global Fourier transform (lines 2395--2409). Every definition, display and label of those lines is retained unchanged. Omitted from this package and not redistributed: 1--441, 449--2076, 2080--2394, 2410--2452, 2534--2858. Nothing inside a retained excerpt is omitted or altered: every retained body is delivered exactly as the article's lines are written, so no omission receipt of any kind accompanies this package. Citations: the retained excerpts carry no citation command at all, so no bibliography entry of the article is transcribed and no reference is redistributed. Reference adaptation: none was needed and none was made; every reference inside the delivered unit resolves inside it, so no unresolved reference or citation is produced. Printed slips kept literal: no printed word, symbol, hypothesis or number of the counted statement or of its proof was altered. Layout and class: the article's own class is \texttt{amsart} with packages including geometry, babel, amsfonts, amsmath, amssymb, amsthm, hyperref, cleveref, bm, tikz-cd, fontenc; the wrapper is the lane's pinned \texttt{amsart} editorial wrapper at 10pt on A4, which restates the theorem-like environments the retained text needs and adds the article's own macro definitions that the retained text needs (\texttt{\textbackslash A}, \texttt{\textbackslash C}, \texttt{\textbackslash Eff}, \texttt{\textbackslash NS}, \texttt{\textbackslash R}, \texttt{\textbackslash Re}, \texttt{\textbackslash adj}), with extra wrapper packages (bm, bbm, dsfont, xspace, cancel, mathtools). The delivered numbering is the unit's own. Assets: no retained span contains a figure environment, a graphics-inclusion command, a PDF-page inclusion, a table or a TikZ picture; no image file is copied into this package, the package declares no asset dependency and no image file is redistributed with it. Licence: version arXiv:2312.13914v2 of this article is distributed under the Creative Commons Attribution 4.0 International licence, as recorded in the arXiv licence metadata for this exact version and shown on the abstract page https://arxiv.org/abs/2312.13914v2. A full rights notice is delivered at licenses/arxiv-2312.13914-CC-BY-4.0.txt. Characters: every retained excerpt is pure ASCII, so the delivered engine, XeLaTeX with its default Latin Modern text font, reports no missing character.
\begin{lemma}\label{Poles of characteristic functions of cones}
	Let $C \subset \Lambda$ be a face, $C_{\R} \subset A_{\R}$ the vector space generated by $C$, $C^{\circ} \subset C$ the interior of $C$ as a subspace of $C_{\R}$ and $C^{\circ}_{\C} := C^{\circ} + i C_{\R}$. Let $b(\ell)$ be the codimension of $C$.
	
	Let $\ell \in C^{\circ}_{\C}$ and $a \in \Lambda^{\circ}$. For each $a \in \Lambda^{\circ}$ the rational function $s^{b(\ell)} \mathcal{X}_{\Lambda}(\ell + as)$ in the complex variable $s$ is holomorphic for $\Re(s) \geq 0$.
	
	If $\ell = 0$ then $s^{\emph{rk} A} \mathcal{X}_{\Lambda}(as) = \mathcal{X}_{\Lambda}(a)$.
\end{lemma}

\begin{equation}\label{eq:log-canonical divisor}
	-K_{(X; \mathbf{A})} = \sum_{D_{\rho} \in D} [D_{\rho}].
\end{equation}

\begin{lemma}\label{lem:Global Fourier transform}
	Let $\chi$ be a quasi-character $\Theta_{\mathbf{A}}(\A_F)/T_{\NS}(\A_F)^1 \to \C^{\times}$. If $\Re(\chi) \in \mathbf{T}_{> (\mathbf{-\frac{3}{2}}, \mathbf{-\frac{5}{2}})}$ then the integral $\hat{\mathcal{G}}(\chi ; \mathbf{s})$ converges absolutely.
	
	If the restriction of $\chi$ to $\mathbf{K}$ or $\mathbf{T}_{\NS}(\A_F)^1$ is non-trivial then $\hat{\mathcal{G}}(\chi ; \mathbf{s}) = 0$.
	
	In general there exist a holomorphic function $\eta_{\mathbf{A}}$ on the tube domain $\mathbf{T}_{> - (\mathbf{\frac{3}{2}}, \mathbf{\frac{5}{2}})}$ such that
	\begin{equation*}
		\hat{\mathcal{G}}(\chi ; \mathbf{s}) = \prod_{\rho \in D / \Gamma} L_{F_\rho}(2; \chi_{\rho}) \eta_{\mathbf{A}}(\chi) \prod_{\substack{v \in \Omega^{\infty} \\ \rho_v \in A_v / \Gamma_v \\ \chi_{\rho_v} = |\cdot|^{z_{\rho_v}}}} \frac{1}{2 + z_{\rho_v}} \mathcal{X}_{\Eff_{(X; \mathbf{A})}}(\mathbf{s} - \mathbf{2} - \bm{z}).
	\end{equation*}
	
	The function $\eta_{\mathbf{A}}$ is rapidly decreasing in vertical strips in the sense that for any compact $C \subset \R_{> -\frac{3}{2}}[D / \Gamma_F] \times \prod_{v \in \Omega_{F}^{\infty}} \R_{> -\frac{5}{2}}[A_v / \Gamma_v]$ and for all $\kappa > 0$ we have for $\chi$ with $\Re(\chi) \in C$ that
	\begin{equation}\label{eq:Rapid decay of eta}
		\eta_{\mathbf{A}}(\chi) \ll_{\kappa, C} (1 + ||\chi||)^{- \kappa}.
	\end{equation}
\end{lemma}

\begin{lemma}\label{lem: Moving the contours}
	There exists an $\varepsilon > 0$ independent of $\chi$ such that the function $s \to (\mathcal{M} \varphi)(-sL )Z_{\mathfrak{a}}(s L, f)$ has a meromorphic continuation to $\Re(s) > a - \varepsilon$. It is holomorphic in this region except for possibly a single pole of order at most $b$ at $s = a$.
	
	For $\bm{z} \in \C[E / \Gamma_F]$, $\bm{w} \in \prod_{v \in \Omega_{F}^{\infty}} \C[B_v/\Gamma_v]$ define $\xi(f, s, \bm{z}, \bm{w})$ as the following sum of integrals
	\begin{equation}\label{eq:Definition of xi}
		c_{\mathbf{A}}^{-1}
		\sum_{\substack{\chi \in (\Theta_{\mathbf{A}}(\A_F)^1/\Theta(F))^{\vee} \\ \chi_{\rho} = 1: \rho \in E/\Gamma_F \\ \chi_{\rho_v} = 1: \rho_v \in B/\Gamma_v}}
		\mathop{\int \cdots \int}_{\substack{\rho \in D_{\adj} / \Gamma_F \\ \Re(z_{\rho}) = 0}}\mathop{\int \cdots \int}_{\substack{v \in \Omega_F^{\infty} \\ \rho_v \in A_{v, \adj} /\Gamma_v \\ \Re(w_{\rho_v}) = 0}} 
		\hat{\mathcal{G}}(\chi \cdot | \cdot|_{\Theta_{\mathbf{A}}}^{(\bm{z}, \bm{w}) + (\bm{z}_{\adj}, \bm{w}_{\adj})}; s L) d\frac{\bm{z}_{\adj}}{2 i \pi} d \frac{\bm{w}_{\adj}}{2 i \pi}.
	\end{equation}
	Where 
	\begin{align*}
		&\bm{z}_{\adj} = (z_{\rho})_{\rho \in D_{\adj} / \Gamma_F} \in \C[D_{\adj} / \Gamma_F], &\bm{w}_{\adj} = (w_{\rho_v})_{\rho_v \in A_{v, \adj}/\Gamma_v} \in \prod_{v \in \Omega_{F}^{\infty}} \C[A_{v, \adj}/\Gamma_v].
	\end{align*}
	
	The sum of integrals $\xi(f, s, \bm{z}, \bm{w})$ converges absolutely for $\Re(s)$ sufficiently large and $\Re(z_{\rho}) > -1, \Re(w_{\rho_v}) > -2$ for $\rho \in E/ \Gamma_F$, $\rho_v \in B_v$ respectively. The function 
	\[
	(s - a)^b \prod_{\rho \in E / \Gamma_F} (z_{\rho} + 1) \prod_{\substack{v \in \Omega_F \\ \rho_v \in B_v/\Gamma_v}} (w_{\rho_v} + 2)\xi(f, s, \bm{z}, \bm{w})
	\] 
	can be analytically continued to the region $\Re(s) > a - \varepsilon, \Re(z_{\rho}) > -1 - \varepsilon, \Re(w_{\rho_v}) > -2 - \varepsilon$ and its value at $s = a, z_{\rho} = -1, w_{\rho_v} = -2 $ is equal to 
	\[
	\lim_{s \to a} (s - a)^b (\mathcal{M} \varphi)(-sL )Z_{\mathfrak{a}}(s L, f).
	\]
	\end{lemma}

\begin{proof}
	Choose $0 < \delta < \frac{1}{4}$ such that 
	\begin{equation}
		K_{(X; \mathbf{A})} + aL - \sum_{\rho \in D_{\adj}} 2\delta [Z_{\rho}] - \sum_{\substack{v \in \Omega_{F}^{\infty} \\ \rho_v \in A_{v, \adj}}} 2\delta [Z_{\rho_v}] \in \Eff_{(X; \mathbf{A})}. 
		\label{Choice of delta}
	\end{equation}
	Such a $\delta$ exists by the definitions of $D_{\adj}$ and $A_{v, \adj}$.
	
	For each $\chi \in (\Theta_{\mathbf{A}}(\A_F)^1/\Theta(F))^{\vee}$ move in the definition of $\mathfrak{Z}(\chi; sL)$ the contours coming from $\rho \in E / \Gamma_F$ and $\rho_v \in B_v / \Gamma_v$ to the lines $\Re(z_{\rho}) = -\frac{5}{4}$ and $\Re(w_{\rho_v}) = - \frac{9}{4}$ respectively. Moreover, also move the contours corresponding to $\rho \in D_{\adj} / \Gamma_F$ and $\rho_v \in A_{v, \adj} / \Gamma_v$ to the lines $\Re(z_{\rho}) = -1 + \delta$ and $\Re(w_{\rho_v}) = - 2 + \delta$ respectively.
	
	By Lemma \ref{lem:Global Fourier transform} each contour moved this way only meets at most a single pole at $z_{\rho} = 1$ and $w_{\rho_v} = 0$ respectively. And it only meets such a pole if $\chi_{\rho_v} = 1$ and $\chi_{\rho_v} = 1$ respectively.
	
	By the residue theorem $\mathfrak{Z}(\chi ; s L)$ is thus equal to the sum over all possible subsets $I \subset E/\Gamma_F$ and $J_v \subset B_v / \Gamma_v$ for all $v \in \Omega_{F}^{\infty}$ such that $\chi_{\rho} = 1$ for all $\rho \in I$ and $\chi_{\rho_v} = 1$ for all $\rho_v \in J_v$ of the following terms
	\begin{equation}\label{eq:Terms after moving the contours}
	\begin{split}
		&\prod_{\rho \in I} \zeta^*_{_{F_{\rho}}}(1) \mathop{\int \cdots \int}_{\substack{ \Re(z_{\rho}) = -\frac{5}{4}: \rho \in I^c \\ \Re(z_{\rho}) = - 1 + \delta: \rho \in D_{\adj}/ \Gamma_{F}}} \prod_{\rho \in I^c \cup (D_{\adj}/ \Gamma_{F})} L(2 + z_{\rho} ; \chi_{\rho}) 
		\mathop{\int \cdots \int}_{\substack{v \in \Omega_{F}^{\infty} \\ 
		\Re(w_{\rho_v}) = -\frac{9}{4}: \rho_v \in J_v^c \\ \Re(w_{\rho_v}) = - 2 + \delta: \rho \in A_{v, \adj} / \Gamma_v}} \\
		&\prod_{\substack{v \in \Omega_{F}^{\infty} \\ \rho_v \in J_v \cup (B_v/ \Gamma_v) \\ \chi_{\rho_v} = 1}} \frac{1}{2 + w_{\rho_v}} 
		\mathcal{X}_{\Eff_{(X; \mathbf{A})}}(s L - \mathbf{2} - (\bm{z}, \bm{w})) \eta_{\mathbf{A}}(|\cdot|^{(\bm{z}, \bm{w})} \cdot \chi)
		d\frac{\bm{z}}{2 i \pi} d \frac{\bm{w}}{2 i \pi}.
	\end{split}
	\end{equation}
	Here $z_{\rho} = 1 $ and $w_{\rho_v} = 1$ for $\rho \in I$ and $\rho_v \in J_v$ respectively.
	
	The integrand decays sufficiently fast in vertical strips for the same reason as for $\mathfrak{Z}(\chi ; \bm{s})$. Similarly, the sum over characters $\chi$ of \eqref{eq:Terms after moving the contours} converges absolutely.
	
	The integrand is a meromorphic function in $s$ and the only part of the integrand which depends on $s$ is the factor $\mathcal{X}_{\Eff_{(X; \mathbf{A})}}(s L - \mathbf{2} - (\bm{z}, \bm{w}))$ so let us analyze this.
	
	By the identity \eqref{eq:log-canonical divisor} $K_{(X ; \mathbf{A})} = (-\mathbf{1}, \mathbf{0})$ we have 
	\begin{equation*}
		\begin{split}
			s L - \mathbf{2} - (\bm{z}, \bm{w}) = &(s - a)L + K_{(X; \mathbf{A})} + aL - 
			\sum_{\rho \in D_{\adj}} 2\delta [Z_{\rho}] - 
			\sum_{\substack{v \in \Omega_{F}^{\infty} \\ \rho_v \in A_{v, \adj}}} 2\delta [Z_{\rho_v}] \; + \\
			&\sum_{\rho \in I^c} (- z_{\rho} - 1) [Z_{\rho}] + \sum_{\rho \in D_{\adj}/ \Gamma_{F}} (- z_{\rho} - 1 + 2 \delta) [Z_{\rho}] \; + \\
			&\sum_{\substack{v \in \Omega_{F}^{\infty} \\ \rho_v \in J_v^{c}}} ( - w_{\rho_v} - 2) [Z_{\rho_v}] + \sum_{\substack{v \in \Omega_{F}^{\infty} \\ \rho_v \in A_{v, \adj} / \Gamma_v}} ( - w_{\rho_v} - 2 + 2 \delta) [Z_{\rho_v}].
		\end{split}
	\end{equation*}
	Let $C_{I, (J_v)_v} \subset \Eff_{(X; \mathbf{A})}$ be the face generated by all $[Z_{\rho}]$ for $\rho \not \in I$ and $[Z_{\rho_v}]$ for $\rho_v \not \in J_v$. We clearly have $C \subset C_{I, (J_v)_v}$. Moreover we have $C = C_{I, (J_v)_v}$ if and only if $I = \emptyset, J_v = \emptyset$ since $[Z_{\rho}], [Z_{\rho_v}] \not \in C$ if $\rho \in E/\Gamma_F, \rho_v \in B_v / \Gamma_v$ by the definition of $E$ and $B_v$ respectively.
	
	Let $b_{I, (J_v)_v}$ be the codimension of $C_{I, (J_v)_v}$. By the above we have $b \geq b_{I, (J_v)_v}$ with equality if and only $I = \emptyset, J_v = \emptyset$
	
 	Note that in the integral \eqref{eq:Terms after moving the contours} we have $ \Re(- z_{\rho} - 1), \Re( z_{\rho} - 1 + 2 \delta), \Re( - w_{\rho_v} - 2), \Re( - w_{\rho_v} - 2 + 2 \delta) > 0$ for $\rho \in I^c, \rho \in E / \Gamma_{\adj}, \rho_v \in J_v^c, \rho_v \in A_{v, \adj} / \Gamma_v$ respectively. 
 	
 	In the notation of Lemma \ref{Poles of characteristic functions of cones} the linear function $s L - \mathbf{2} - (\bm{z}, \bm{w})$ is thus equal to a sum of $(s - a)L$ and an element in $C_{I, (J_v)_v, \C}^{\circ}$. The lemma then implies that there exists a $\varepsilon > 0$ such $\mathcal{X}_{\Eff_{(X; \mathbf{A})}}(s L - \mathbf{2} - (\bm{z}, \bm{w}))$ has a meromorphic continuation to the region $\Re(s) > a - \varepsilon$ which is holomorphic except for a single pole at $s = a$ of order at most $b_{I, (J_v)_v}$. It follows that the same is true for the integral \eqref{eq:Terms after moving the contours}. 
 	
 	The integral $\xi(f, s, \bm{z}, \bm{w})$ converges and has a meromorphic continuation for the same reason.
 	
 	 By moving the contours corresponding to $z_{\rho}$ for $\rho \in D_{\adj}/\Gamma_F$ and $w_{\rho_v}$ for $\rho_v \in A_{v, \adj}/ \Gamma_v$ back we see that the contribution of \eqref{eq:Terms after moving the contours} for the case $I = J_v = \emptyset$ is equal to the value at $z_{\rho} = -1$ for $\rho \in E/\Gamma_v$ and $w_{\rho_v} = -2$ for $\rho_v \in B_v/ \Gamma_v$ of the meromorphic continuation of
 	\[
 	\prod_{\rho \in E / \Gamma_F} (z_{\rho} + 1) \prod_{\substack{v \in \Omega_F \\ \rho_v \in B_v/\Gamma_v}} (w_{\rho_v} + 2)\mathop{\int \cdots \int}_{\substack{\rho \in D_{\adj} / \Gamma_F \\ \Re(z_{\rho}) = 0}}\mathop{\int \cdots \int}_{\substack{v \in \Omega_F^{\infty} \\ \rho_v \in A_{v, \adj} /\Gamma_v \\ \Re(w_{\rho_v}) = 0}} 
 	\hat{\mathcal{G}}(\chi \cdot | \cdot|_{\Theta_{\mathbf{A}}}^{(\bm{z}, \bm{w}) + (\bm{z}_{\adj}, \bm{w}_{\adj})}; s L) d\frac{\bm{z}_{\adj}}{2 i \pi} d \frac{\bm{w}_{\adj}}{2 i \pi}.
 	\] 
 	
 	We can switch the limit with the sum over characters by absolute convergence.
\end{proof}

\end{document}
